\section{Results}
\label{sec:results}

\begin{table}[t]
\centering
\small
\caption{Main split-scoped results. Rows are evidence for the stated slices, not universal benchmark claims. Fix/harm/wrong-forward is reported where available; a dash denotes an audit-clean aggregate without per-item switch accounting in the table source.}
\label{tab:main-results}
\begin{tabular}{lcccll}
\toprule
Slice & ARG & PromptNode & $\Delta$ & Fix/Harm & Claim role \\
\midrule
AQuA\_audit\_clean\_tail31 & 28/31 & 31/31 & 3 & - & positive\_clean\_improvement \\
AQuA\_conservative\_remove\_114\_tail30 & 28/30 & 30/30 & 2 & - & positive\_under\_boundary\_label\_removal \\
MMLU\_STEM\_main97 & 88/97 & 95/97 & 7 & 7/0 & positive\_with\_completed\_controls \\
MMLU\_nonSTEM40\_safety\_replay & 37/40 & 37/40 & 0 & 0/0 & no\_harm\_boundary\_only \\
MMLU\_nonSTEM40\_raw\_warning\_override & 37/40 & 34/40 & -3 & 0/3 & negative\_transfer\_limitation \\
MMLU\_DefinitionChallenge\_fresh24 & 21/24 & 22/24 & 1 & 1/0/0 & targeted\_nonmath\_positive\_no\_harm \\
\bottomrule
\end{tabular}
\end{table}

\begin{table}[t]
\centering
\small
\caption{Cost evidence rows. Direct or existing-summary costs are separated from the MMLU STEM probe estimate, which is not exact original R306 accounting.}
\label{tab:cost-results}
\begin{tabular}{lcccl}
\toprule
Row & Score & Cost & Cost/task & Evidence \\
\midrule
AQuA\_ARG\_baseline & 28/34 & 0.93285 & 0.027437 & direct\_or\_existing\_summary \\
AQuA\_R068\_cost\_reduced & 31/34 & 3.597409 & 0.105806 & direct\_or\_existing\_summary \\
AQuA\_R267\_safety\_cost & 31/31 & 3.597409 & 0.116045 & direct\_or\_existing\_summary \\
MMLU\_STEM\_R306\_probe\_estimate & 95/97 & 0.676454 & 0.006974 & probe\_estimate\_not\_exact \\
MMLU\_DefinitionChallenge\_R413 & 22/24 & 0.220095 & 0.009171 & direct\_or\_existing\_summary \\
MMLU\_plain\_requery\_R316 & 37/40 & 0.11794 & 0.002949 & direct\_or\_existing\_summary \\
\bottomrule
\end{tabular}
\end{table}


Table~\ref{tab:main-results} reports split-scoped accuracy evidence rather than universal benchmark claims. Table~\ref{tab:cost-results} records cost provenance and separates direct or existing-summary costs from probe estimates.

\paragraph{AQuA audited tail.}
On the audit-clean AQuA tail31 slice, ARG obtains 28/31 while the PromptNode graph policy obtains 31/31. Under the conservative boundary that also removes task 114, ARG obtains 28/30 and the PromptNode variant obtains 30/30. These results support a scoped AQuA improvement claim, not a claim about all AQuA variants or noisy labels.

\paragraph{MMLU STEM transfer.}
On the MMLU STEM main97 slice, the PromptNode stack obtains 95/97 compared with 88/97 for ARG, with fix/harm 7/0 in the evidence table. The draft cost row is based on a 12-task cost probe: R425 estimates a combined main97 cost of 0.676454 for ARG direct plus PromptNode-stack style calls. This cost is marked as a probe estimate and must not be described as exact R306 accounting.

\paragraph{Non-STEM transfer.}
The broader mixed non-STEM result supports only a no-harm boundary: R396 reports ARG 41/48 and the retrained hybrid 41/48, with no harms but no positive delta. A targeted fresh DefinitionChallenge run gives a narrower positive result: R413/R414 reports ARG 21/24 versus DefinitionChallenge 22/24, with fix/harm/wrong-forward 1/0/0. The raw candidate pool alone is 21/24, so the improvement comes from graph selection rather than from a stronger plain requery candidate pool.

\paragraph{Controls and limitations of evidence.}
Plain requery is not enough to explain the positive targeted result: R316 plain requery on non-STEM40 has the same 37/40 accuracy as ARG and fix/harm 1/1. However, this does not yet establish broad transfer. The DefinitionChallenge result is a targeted nonmath result, the lifecycle trace is replayed from R413 outputs, and the external integrity audit is pending.
